Searches for dense packings of congruent shapes move rigid pieces. Hardening starts each piece as its inscribed disk and rounds it into the real shape while the container is compressed. We ask whether it is a search worth running alongside existing ones, and test it at matched budget and comparable tuning against rigid starts, growth, simulated annealing and perturbation--compression on seven families of squares, triangles, hexagons and disks in squares, triangles and circles (\CsevenInstances{} instances), scoring each method by the lowest container size it reaches. Hardening is not a better default: tuned rigid starts reach the lowest size most often. It is a different search. On squares in a square it reaches packings the other four methods miss, including Trump's $n=11$ and Wainwright's $n=19$ in 32 runs per method, and in a later confirmatory campaign at ten times the budget (16 runs per method) the only one to reach the best known packings for $n=\LongHardenOnlyReachedNewNs$, where annealing alone reaches $n=\LongSaOnlyReachedNs$. In exploratory ablations, compressing disks and then switching to polygons at once, or growing rigid polygons along hardening's area schedule, shows no detectable difference from rigid starts, and hardening's unique packings contain more tilted pieces. The bias costs elsewhere: hardening is worse on triangles and on four families held out from all development, where a short pilot that chooses the path per instance does no better than always starting rigid. Shape continuation is an alternative search that reaches different packings, worth keeping beside rigid ones; we did not show that splitting a fixed budget between them pays, and we have no test that predicts where it helps. Every experiment was planned before it ran, every run is replayable, and every number is generated from an event log.
